\documentclass[a4paper, 12pt]{article}
\usepackage[obeyspaces]{url}

%\usepackage{etoolbox}
\usepackage[utf8]{inputenc}
\usepackage[spanish, es-tabla]{babel}
\usepackage{graphicx}
\usepackage{pdfpages}
\usepackage{xcolor}
\definecolor{LightGray}{gray}{0.9}
\usepackage{hyperref}
\usepackage[caption = false]{subfig}
\usepackage{pythontex}
\usepackage{float}
\usepackage{nicematrix} % To get nice tables
\usepackage{amssymb}
\usepackage{cancel} % Cancel terms in an equation
\usepackage{rsfso} % Reluctance
\usepackage{bm}
\usepackage{xcolor}
\usepackage{pifont} % more styles for bullets
\decimalpoint
\hypersetup{
	colorlinks=true,
	linkcolor=blue,
	urlcolor=blue
	%pdfpagemode=FullScreen
	}
\urlstyle{same}
\setlength{\parskip}{1.3em}
\usepackage[top=2.5cm,bottom=2.5cm,left=2.5cm,right=2.5cm,headsep=8pt,a4paper]{geometry}



\begin{document}
\begin{titlepage}
     \newcommand{\HRule}{\rule{\linewidth}{0.5mm}}
     \center
     \textsc{\LARGE  Universidad Nacional de San \\[0.2cm] Antonio Abad del Cusco}\\[1.5cm]
	 \includegraphics[scale=0.3]{figures/escudo.jpg}\\[1cm]
     \textsc{\Large Facultad de Ingeniería Eléctrica, \\Electrónica, Informática y Mecánica}\\[0.5cm]
     \textsc{\large Escuela Profesional de Ingeniería Electrónica}\\[5mm]
	 
	 \HRule \\[0.4cm]
     { \huge \bfseries CIRCUITOS ELECTRÓNICOS I}\\[1.8mm] 
	 \HRule\\[5mm]
     \begin{minipage}{\textwidth}
        \center
        \textit{Trabajo final}\\[0.2cm]
        \text{\large Diseño de un amplificador}\\[3mm]
		\begin{tabular}{ll}
        \emph{Docente:} \\[1.8mm]
            \hspace{1.4cm}\textsc{Ing. Milton Jhon Velasquez Curo}\\[3mm]
        \emph{Estudiantes:} & \quad \emph{Códigos:}\\[1.8mm]
            \hspace{1.4cm} \textsc{Roly Sandro Gutierrez Benito} & \hspace{1.7cm} 182967\\[3mm]
            \hspace{1.4cm} \textsc{Edison Abado Ancco} & \hspace{1.7cm} 145012\\[3mm]
            \hspace{1.4cm} \textsc{Alberto Arturo Quiñones Pacco}& \hspace{1.7cm} 160901\\[3mm]
		 \end{tabular} \\[3.5mm]
		 \textsc{CUSCO-PERÚ}\\
         \textsc{2022}
     \end{minipage}
\end{titlepage}
\newpage
\tableofcontents

\newpage
\section{Diseñar e implementar un amplificador de señales}
Las características del amplificador de señales son los siguientes:
\begin{itemize}
    \item Señal de entrada pico $V_{ip} = 0.1 - 1\ \text{mV}$
    \item Ganancia de voltaje $A_v = 80\ \text{dB}$
    \item Impedancia de salida $Z_o = 8\ \Omega$
    \item Rango de frecuencias $f_l = 40 \text{Hz}$, $f_h = 40\ \text{kHz}$ 
\end{itemize}

\section{Etapas del amplificador}

\subsection{Preamplificador}
Para esta etapa, se hizo uso de un amplificador multietapa.
\begin{itemize}
    \item \textbf{Primera etapa}
        \begin{itemize}
            \item Polarización universal en emisor común.
            \item Transistor \verb|BC548|
        \end{itemize}

    \item \textbf{Segunda etapa}
        \begin{itemize}
            \item Polarización universal en emisor común.
            \item Transistor \verb|BC548|
        \end{itemize}
\end{itemize}

\subsection{Diseño y cálculos para la primera etapa}
Dado el siguiente circuito en \textbf{polarización universal} en \textbf{emisor común}:  
\begin{figure}[H]
    \centering
    \includegraphics[scale=0.5]{figures/image1.png}
    \caption{Primera etapa del preamplificador con polarización universal EC}
    \label{fig:image1}
\end{figure}
\subsubsection{Gráficas del punto de operación}
\subsubsection{Rectas de carga en DC y AC}
\subsubsection{Gráfica de Bode}
\subsubsection{Potencia}

\subsection{Diseño y cálculos para la segunda etapa}
\subsubsection{Gráficas del punto de operación}
\subsubsection{Rectas de carga en DC y AC}
\subsubsection{Gráfica de Bode}
\subsubsection{Potencia}

\newpage
\subsection{Amplificador de potencia}
En esta etapa, se tienen las siguientes características: 
\begin{itemize}
    \item Amplificador clase B
\end{itemize}

\subsubsection{Diseño y cálculos}
\subsubsection{Gráficas del punto de operación}
\subsubsection{Rectas de carga en DC y AC}
\subsubsection{Gráfica de Bode}
\subsubsection{Potencia}
\end{document}
